\section{Modularity and Reconfigurability}
\label{sec:modularity}

If Section~\ref{sec:shape-changing} ended with shape as the form of an integral body, this section asks what changes when shape is no longer integral but assembled. The literature of modular reconfigurable robotics (MRR) treats the body as a collection of comparable units whose arrangement can be changed, by design or in operation, to suit different tasks.

\subsection{Concept}
\label{subsec:mrr-Concept}
 

An MRR system, in the formulation given by Liang et al.'s recent review (\textit{\cref{fig:liang_mrr_timeline}}), is a collection of independent modules that can be physically connected, with overall morphology adjusted by changing which modules are connected and how \cite{liangModularReconfigurableRobots2026}. Three advantages are commonly attributed to this design philosophy. \textit{Versatility}, since the same modules support different morphologies for different tasks. \textit{Robustness}, since damage to individual modules can be compensated by reconfiguration. \textit{Cost-effectiveness}, since a small set of standardized modules amortizes across many possible robots. The most expansive vision in this tradition is the ``bucket of stuff'' image: a loose collection of modules that, given a task, autonomously assemble into the morphology the task requires \cite{liangModularReconfigurableRobots2026}.

Liang and colleagues structure the field along three application domains, namely locomotion, manipulation, and construction, and along five increasing levels of reconfigurability, from basic chain-and-lattice rearrangement at Level 1 to dynamic flow morphing in unstructured environments at Level 5 \cite{liangModularReconfigurableRobots2026}. Within locomotion, two regimes are distinguished: \textit{gait locomotion} keeps a fixed configuration during motion and moves through coordinated module action, while \textit{flow locomotion} dynamically reconfigures the body in response to its environment. Gait locomotion is primarily associated with the lower reconfigurability levels, where the configuration is decided at assembly and held during operation; flow locomotion belongs to the higher levels, where modules can re-attach and re-arrange while the robot is in motion.

What distinguishes MRR from the shape-changing systems of Section~\ref{sec:shape-changing} is the centrality of the \textit{connection}. In an integral shape-changing robot, the form is decided at the design stage and welded to the device's identity; in an MRR, the form emerges from how repeatable units have been joined. The choice of connection mechanism, including its degrees of freedom, strength, area of acceptance, and directionality, therefore defines what configurations a given module set can take, and how easily one configuration can become another. Connection design, in other words, is the design space of MRR.

\subsection{Related Works}
\label{subsec:mrr-related}

A first thread of recent work treats connection topology itself as the contribution. Tu, Liang, and Lam's FreeSN \cite{tuFreeSNFreeformStrutnode2022} introduces a strut-node architecture in which spherical-shell nodes accept strut modules through freeform magnetic connectors. The authors argue that earlier MRR designs, most of them based on cubic modules with at most six fixed-point connectors, are constrained by a topological deficiency: with a minimum central angle of 90 degrees between connectors, cubic-module configurations cannot form triangles, and the resulting structures lack stability at scale. FreeSN's spherical nodes lower this minimum angle and enable triangulated substructures, expanding the topological space available to the assembled robot. Chin et al.'s AuxBot \cite{chinFlipperStyleLocomotionStrong2023}, already encountered for its shape-changing modules, contributes on a different axis: between modules, flexible wire constraints replace rigid connectors, allowing volume-changing modules to expand past one another and producing the flipper-style locomotion observed in mudskippers and sea turtles (\textit{\cref{fig:auxbot}}). Yu, Matthews, Wang, and colleagues' autonomous modular legs \cite{AgileLeggedLocomotion} extend the connection idea further, treating each leg as a fully self-contained agent that can be bolted to any other; the body that emerges depends entirely on which legs have been joined where, and the authors encode this body design space into a latent representation that allows automatic discovery of novel multi-legged forms.

A second thread approaches the same question from the opposite end of the connection spectrum. Li et al.'s particle robotics \cite{liParticleRoboticsBased2019} replaces deliberate connection with stochastic coupling: many disk-shaped particles, each capable only of volumetric oscillation, are loosely held together by dangling magnets. Collective locomotion, object transport, and phototaxis emerge from this loosely-coupled aggregate, and the system continues to function even when twenty percent of particles fail. Particle robotics is a different paradigm from configured MRR, but it points to the same lesson from the opposite direction: when many simple modules interact, aggregate behavior depends on how they are arranged and coupled, not only on what each module does on its own. Alongside these, a small body of recent work begins to push MRR toward soft and electrohydraulic actuation. Zhao et al.'s shape-changing tensegrity-blocks \cite{zhao_modular_2025} position themselves explicitly against rigid-module mainstream MRR, combining tensegrity structures with lightweight deformable units; and Yoder et al.'s HEXEL system \cite{yoderHexagonalElectrohydraulicModules2024}, returning here from Section~\ref{sec:shape-changing}, demonstrates that modular bodies built from soft electrohydraulic actuators can achieve the high-speed untethered operation that pneumatic-driven soft modular systems have so far struggled with.

\subsection{Gap}
\label{subsec:mrr-gap}

The literature surveyed above is rich in connection mechanisms, design space frameworks, and reconfiguration strategies. Two of its underdeveloped corners, however, are exactly the corners this thesis works in.

The first is material. Most MRR systems remain rigid-bodied. The systematic exploration of compliant material as a computational substrate, established as a prospect in Section~\ref{sec:soft-robotics}, has barely entered the modular literature. A small but growing set of works has begun to push in this direction, including tensegrity-blocks, electrohydraulic modules, and inflatable trusses, but the field is still organized around the assumption that each module is, individually, a rigid object. The implications of building MRR from compliant, passively coupled modules, where neighbor inflation deforms a module's effective geometry rather than merely connecting to it, have not been systematically explored.

The second is actuation. Motor-driven actuation dominates the modular literature, both for kinematic clarity and for power density. Pneumatic actuation, despite its appeal for soft-body designs as established in Section~\ref{sec:soft-robotics}, remains rare in MRR; Zhao and colleagues note explicitly that untethered pneumatic operation is still an open challenge for soft modular systems. The intersection of \textit{soft}, \textit{modular}, and \textit{pneumatic} is therefore a thin region of the design space, occupied by a handful of research prototypes and not yet supported by systematic comparison across configurations.

These two gaps inherit, and sharpen, the gap left by Section~\ref{sec:shape-changing}. There, the missing element was comparison across geometries; here, even within the still-emerging soft-modular line, no comparative study yet asks how different geometric arrangements of identical compliant modules shape what the robot can do. The next section turns to one further dimension of the problem: given that body and configuration can take on part of the work, how should the controller distribute the rest? It examines emergence and bottom-up coordination as a complementary paradigm.
